% Generated by scripts/assemble.py; corrections belong in source/corrections.json.
\examhead{27}{中国矿业大学(北京) 2026 年数学分析试卷(回忆版)}{31}
% p032-b01; source PDF page 32
\blockanchor{p032-b01}
(注.题目来自考生回忆,部分表述或数据可能有误)\par
% p032-b02; source PDF page 32
\par\Needspace{5\baselineskip}
\blockanchor{p032-b02}
\majorquestion{一}填空题.\par
% p032-b03; source PDF page 32
\blockanchor{p032-b03}
\question{1}$\lim_{\alpha \to 0}\int_{\alpha}^{1 + \alpha}\frac{\mathrm{d}x}{1 + x^2 + \alpha^2} =$\par
% p032-b04; source PDF page 32
\blockanchor{p032-b04}
\question{2}$\sum_{n=1}^{\infty} (-1)^{n} \frac{2n - 1}{(2n + 3)!} =$ \_\_\_\_.\par
% p032-b05; source PDF page 32
\blockanchor{p032-b05}
\question{3}$\int \frac{\ln x}{(1 - x)^2} \, \mathrm{d}x =$\par
% p032-b06; source PDF page 32
\blockanchor{p032-b06}
\question{4}高阶无穷小问题\par
% p032-b07; source PDF page 32
\blockanchor{p032-b07}
\question{5}$f(x) = \frac{1}{x^2}\arctan \frac{1}{x - \frac{1}{x}}$ 的第一类间断点\par
% p032-b08; source PDF page 32
\blockanchor{p032-b08}
\question{6}$x = a\cos t, y = a\sin t, z = at,$ 求曲线积分 $\int_{L}\frac{z^{2}}{x^{2} + y^{2}}\mathrm{d}s =$\par
% p032-b09; source PDF page 32
\blockanchor{p032-b09}
\question{7}$\int_0^{2\pi}\mathrm{d}x\int_x^{2\pi}\frac{\sin^2y}{y^2}\mathrm{d}x =$\par
% p032-b10; source PDF page 32
\blockanchor{p032-b10}
\question{8}反常积分  $\int_{0}^{+\infty}\frac{\sin x}{x}dx=$  \_\_\_\_ (填“绝对收敛”或“条件收敛”或“发散”).\par
% p032-b11; source PDF page 32
\blockanchor{p032-b11}
\question{9}$y = f(x)$ 与 $y = x^{2} - x$ 在 $(1,0)$ 点有公共的切线，求 $\lim_{n\to \infty}f\left(\frac{n}{n + 2}\right) = \underline{\quad}$.\par
% p032-b12; source PDF page 32
\blockanchor{p032-b12}
\question{10}求 $y = ax, y = bx, x + y = c, x + y = d$ 所围曲面的面积\par
% p032-b13; source PDF page 32
\par\Needspace{5\baselineskip}
\blockanchor{p032-b13}
\majorquestion{二}证明 $f(x) = e^{x} \cos \frac{1}{x}$ 在 (0,1) 上的一致连续性.\par
% p032-b14; source PDF page 32
\par\Needspace{5\baselineskip}
\blockanchor{p032-b14}
\majorquestion{三}求 $f(x) = \frac{2x^2}{(1 - x)^2}$ 的极值点与极值, 凹凸区间, 拐点与渐近线.\par
% p032-b15; source PDF page 32
\par\Needspace{5\baselineskip}
\par\Needspace{6\baselineskip}
\blockanchor{p032-b15}
\majorquestion{四}求曲面积分\par
\nopagebreak[4]
% p032-b16; source PDF page 32
\blockanchor{p032-b16}
\displayformula{\iint \frac {x ^ {2} y \mathrm{d} y \mathrm{d} z + y ^ {2} z \mathrm{d} z \mathrm{d} x + (z ^ {2} + 1) \mathrm{d} x \mathrm{d} y}{4 x ^ {2} + 4 y ^ {2} + z ^ {2}}.}
% p032-b17; source PDF page 32
\blockanchor{p032-b17}
其中  $\Sigma: 4x^{2} + 4y^{2} + z^{2} = 4, z \geq 0.$\par
% p032-b18; source PDF page 32
\par\Needspace{5\baselineskip}
\blockanchor{p032-b18}
\majorquestion{五}叙述有限覆盖定理与连续函数的根的存在性定理, 并用有限覆盖定理证明根的存在性定理.
\majorquestion{六}证明  $f_{n}(x)=nx(1-x)^{n}$  在  $[0,1]$  上的一致收敛性.\par
% p032-b19; source PDF page 32
\par\Needspace{5\baselineskip}
\blockanchor{p032-b19}
\majorquestion{七}设 $a_1 = 3, a_{n+1} = \frac{1}{2} a_n + \frac{1}{a_n}$.\par
% p032-b20; source PDF page 32
\blockanchor{p032-b20}
(1) 证明: $\{a_{n}\}$ 收敛, 并求其极限.\par
% p032-b21; source PDF page 32
\blockanchor{p032-b21}
(2) 证明: 对任意正数 $p$, 数项级数 $\sum_{n=1}^{\infty} n^{p} \left( \frac{a_{n}}{a_{n+1}} - 1 \right)$ 收敛.\par
% p032-b22; source PDF page 32
\par\Needspace{5\baselineskip}
\par\Needspace{6\baselineskip}
\blockanchor{p032-b22}
\majorquestion{八}$f(x)$ 在 $[0,1]$ 上连续, $I = \int_0^1 f(x)\mathrm{d}x\neq 0$ ，证明：在 $[0,1]$ 内存在两个不同的点 $x_{1},x_{2}$ ，使得\par
\nopagebreak[4]
% p032-b23; source PDF page 32
\blockanchor{p032-b23}
\displayformula{\frac {1}{f (x _ {1})} + \frac {1}{f (x _ {2})} = \frac {2}{I}.}
% p032-b24; source PDF page 32
\par\Needspace{5\baselineskip}
\blockanchor{p032-b24}
\majorquestion{九}求 $f(x, y) = 3(x - 2y)^2 + x^3 - 8y^3$ 的极值.\par
